\begin{definition}[Representative word-tuple separation]
    \label{def:representative_word_tuple_span}
    \lean{MPSTensor.SectorDecomposition.RepresentativeWordTupleSpanAt}
    \lean{MPSTensor.SectorDecomposition.EventuallyRepresentativeWordTupleSpan}
    \leanok
    \uses{def:sector_weight_data, rem:word_tuple_span_top}
    Let $(A_j)_{j=0}^{g-1}$ be the basis representatives of a sector
    decomposition. They have \emph{representative word-tuple separation at
    length $L$} when
    \begin{align}
        \spn_{\C}\bigl\{(A_j^w)_{j=0}^{g-1}:
        w\in\{0,\ldots,d{-}1\}^{L}\bigr\}
        &=\prod_{j=0}^{g-1}\MN{D_j}.
        \label{eq:mpdo_rep_tuple_sep}
    \end{align}
    They have \emph{eventual representative word-tuple separation} when this
    equality holds at every sufficiently large length. Notice that the
    family is indexed by the minimal representatives $j$, and not by the
    repeated copies lying over each $j$.
\end{definition}

\begin{theorem}[Representative-grouped Lemma L]
    \label{thm:representative_grouped_insert_eq}
    \lean{MPSTensor.SectorDecomposition.insertedTensor_basis_eq_of_firstSiteActionAgree_of_coeff_ne_zero}
    \lean{MPSTensor.SectorDecomposition.insertedTensor_basis_eq_of_firstSiteActionAgree}
    \lean{MPSTensor.SectorDecomposition.insertedTensor_basis_eq_of_sameMPV₂Pos_firstSiteActionAgree}
    \leanok
    \uses{def:representative_word_tuple_span,
      def:sector_weight_data, cor:nonvanishing_finite_power_sum}
    Let $A_{\mathrm{tot}}$ be the tensor of a sector decomposition with basis
    representatives $(A_j)_{j=0}^{g-1}$, multiplicities $r_j>0$, nonzero weights
    $\mu_{j,q}$, and grouped coefficients
    $c_j^{(N)}=\sum_{q=0}^{r_j-1}\mu_{j,q}^{N}$.
    Suppose that the first-site actions of $Y,Z\in\MN{d}$ agree on the MPV
    family of $A_{\mathrm{tot}}$. If representative word-tuple separation
    holds at length $L$ and $c_{j_0}^{(L+1)}\neq0$, then
    $YA_{j_0}=ZA_{j_0}$.
    Consequently, under eventual representative word-tuple separation,
    $YA_j=ZA_j$ for every representative $j$. The same conclusion holds if
    the first-site equality is initially given for any tensor with the same
    positive-length MPV family as $A_{\mathrm{tot}}$.

    The representative grouping is the argument of
    \cite[Appendix~C.3, lines~1835--1858]{Cirac2016MPDO_arXiv}.

    In the first auxiliary assertion, eventual representative word-tuple
    separation is an explicit assumption. The later post-blocked theorem
    derives the needed finite representative separation from the full BNT
    canonical-form hypotheses. The sharp $3D^5$ representative-separation
    bound is Theorem~\ref{thm:sharp_bnt_block_separation}; its detailed source
    comparison is recorded in
    \path{docs/paper-gaps/cpgsv17_bicf_block_separation.tex}.
\end{theorem}

\begin{proof}
    \leanok
    \uses{def:representative_word_tuple_span,
      lem:paper_decomp_formula, cor:nonvanishing_finite_power_sum}
    At a fixed tail length $L$, expand the first-site equality at system size
    $L+1$ and group all copies lying over the same representative. For every
    length-$L$ word $w$ this gives
    \begin{align}
        \sum_{j=0}^{g-1}
        \tr\!\left(c_j^{(L+1)}(YA_j-ZA_j)A_j^w\right)
        &=0.
        \notag
    \end{align}
    Representative word-tuple separation therefore forces each matrix
    $c_j^{(L+1)}(YA_j-ZA_j)$ to vanish. A nonzero coefficient may then be
    cancelled.

    For the eventual assertion, choose a threshold $L_0$ and put
    $M=L_0+1$. For a fixed representative $j$, apply
    Corollary~\ref{cor:nonvanishing_finite_power_sum} to the nonzero numbers
    $\mu_{j,q}^{M}$. It yields $r\geq1$ such that
    $c_j^{(Mr)}=\sum_q(\mu_{j,q}^{M})^r\neq0$.
    Since $Mr-1\geq L_0$, the fixed-length argument applies with
    $L=Mr-1$. Finally, equality of the positive-length MPV families gives the
    first-site hypothesis for the sector-decomposition tensor.
\end{proof}

\begin{definition}[Direct sum with one marked letter]
    \label{def:marked_representative_direct_sum}
    \lean{MPSTensor.SectorDecomposition.markedTensor}
    \leanok
    \uses{def:sector_weight_data}
    Let
    \begin{align}
        A_{\mathrm{tot}}^i
        &=\bigoplus_{j=0}^{g-1}\bigoplus_{q=0}^{r_j-1}
          \mu_{j,q}A_j^i,
        \label{eq:mpdo_sector_direct_sum}
    \end{align}
    be the tensor of a sector decomposition. Let $s$ range over a second
    finite alphabet, and for each representative $A_j$ let $C_j^s$ be a
    matrix on its bond space. The corresponding direct sum with one marked
    letter is
    \begin{align}
        C_{\mathrm{tot}}^s
        &=\bigoplus_{j=0}^{g-1}\bigoplus_{q=0}^{r_j-1}
          \mu_{j,q}C_j^s.
        \label{eq:mpdo_marked_direct_sum}
    \end{align}
    The marked alphabet need not coincide with the alphabet of the unmarked
    tensor. This is the algebraic form of the marked matrices in the two
    diagrams of
    \cite[Proposition~4.13, lines~1909--1919]{Cirac2016MPDO_arXiv}.
\end{definition}

\begin{lemma}[Trace expansion with one marked letter]
    \label{lem:marked_representative_trace_expansion}
    \lean{MPSTensor.SectorDecomposition.trace_markedTensor_mul_evalWord}
    \leanok
    \uses{def:marked_representative_direct_sum, def:sector_weight_data}
    For every marked index $s$ and every unmarked word $w$,
    \begin{align}
        \tr\!\left(C_{\mathrm{tot}}^s A_{\mathrm{tot}}^w\right)
        &=\sum_{j=0}^{g-1}c_j^{(|w|+1)}
          \tr\!\left(C_j^s A_j^w\right),
        \label{eq:mpdo_marked_trace_expansion}\\
        c_j^{(N)}
        &=\sum_{q=0}^{r_j-1}\mu_{j,q}^N.
        \notag
    \end{align}
    The exponent $|w|+1$ includes the weight at the marked site.
\end{lemma}

\begin{proof}
    \leanok
    Expand both matrices as direct sums. Their product is the direct sum of
    $\mu_{j,q}^{|w|+1}C_j^sA_j^w$.
    Taking the trace and summing first over $q$ gives the stated coefficient
    $c_j^{(|w|+1)}$.
\end{proof}

\begin{theorem}[Algebraic marked-letter extension of representative separation]
    \label{thm:representative_grouped_marked_separation}
    \lean{MPSTensor.SectorDecomposition.markedTensor_basis_eq_of_trace_agree_of_coeff_ne_zero}
    \lean{MPSTensor.SectorDecomposition.markedTensor_basis_eq_of_trace_agree}
    \leanok
    \uses{def:representative_word_tuple_span,
        def:marked_representative_direct_sum,
        lem:marked_representative_trace_expansion,
        cor:nonvanishing_finite_power_sum}
    Let $(C_j^s)$ and $(E_j^s)$ be two marked families over the same sector
    decomposition. Suppose that, for every marked index $s$ and every
    unmarked word $w$,
    \begin{align}
        \tr\!\left(C_{\mathrm{tot}}^s A_{\mathrm{tot}}^w\right)
        &=\tr\!\left(E_{\mathrm{tot}}^s A_{\mathrm{tot}}^w\right).
        \notag
    \end{align}
    If representative word-tuple separation holds at length $L$ and
    $c_{j_0}^{(L+1)}\neq0$, then $C_{j_0}=E_{j_0}$. Under eventual
    representative word-tuple separation, one has $C_j=E_j$ for every $j$.

    Lemma~L is stated for marks induced by physical operators on the first
    site. After its closed-chain identity has been expanded, the separation
    calculation uses only the matrices in the marked slot. The theorem is
    the resulting algebraic extension of
    \cite[Appendix~C.3, lines~1848--1858]{Cirac2016MPDO_arXiv}; it is not
    attributed to the paper as a separately stated arbitrary-mark theorem.
\end{theorem}

\begin{proof}
    \leanok
    \uses{def:representative_word_tuple_span,
        lem:marked_representative_trace_expansion,
        cor:nonvanishing_finite_power_sum}
    Fix $s$ and a tail length $L$, and put
    $\Delta_j=c_j^{(L+1)}(C_j^s-E_j^s)$.
    Lemma~\ref{lem:marked_representative_trace_expansion} gives
    $\sum_{j=0}^{g-1}\tr(\Delta_jA_j^w)=0$ for every length-$L$ word $w$.
    Representative word-tuple separation implies $\Delta_j=0$ for every
    $j$. If $c_{j_0}^{(L+1)}\neq0$, cancellation gives
    $C_{j_0}^s=E_{j_0}^s$. Since $s$ was arbitrary,
    $C_{j_0}=E_{j_0}$.

    Under eventual separation, apply
    Corollary~\ref{cor:nonvanishing_finite_power_sum} to the nonzero weights
    above each fixed representative, exactly as in
    Theorem~\ref{thm:representative_grouped_insert_eq}. This selects a
    sufficiently large $L$ with $c_j^{(L+1)}\neq0$ and proves $C_j=E_j$.
\end{proof}

% Source anchors: CPSV16 Appendix C.3, lines 1835--1858 and 1909--1919.
\begin{theorem}[Full-coordinate marked trace identity]
    \label{thm:cpsv_full_coordinate_marked_trace_identity}
    \lean{MPSTensor.CPSVCanonicalFormData.ActiveBNTRefinement.%
        groupedMarkedBlocks}
    \lean{MPSTensor.CPSVCanonicalFormData.ActiveBNTRefinement.%
        groupedMarkedTensor}
    \lean{MPSTensor.CPSVCanonicalFormData.ActiveBNTRefinement.%
        trace_groupedMarkedTensor_eq_representative_markedTensor}
    \leanok
    \uses{thm:cpsv_canonical_form_active_bnt_refinement,
        thm:cpsv_active_representative_sector_decomposition,
        def:marked_representative_direct_sum,
        lem:marked_representative_trace_expansion}
    Let $T_{\mathrm{grp}}$ be the full grouped tensor of
    Theorem~\ref{thm:cpsv_canonical_form_active_bnt_refinement}, and let
    $P_{\mathrm{act}}$ be its active representative sector decomposition.
    For a marked family $(C_j^s)_j$, define the marked tensor
    $\widehat C_{\mathrm{grp}}$ on all grouped listed coordinates as follows.
    The active copy $(j,q)$ carries the marked block
    $\lambda_{j,q}C_j^s$, where
    $\lambda_{j,q}=\nu_{k(j,q)}\zeta_{k(j,q)}$, and every inactive listed
    coordinate carries the zero block. Thus all original listed coordinates
    remain present.

    For every marked letter $s$ and every unmarked word $w$, one has
    \begin{align}
        \tr\!\left(\widehat C_{\mathrm{grp}}^sT_{\mathrm{grp}}^w\right)
        &=\tr\!\left(C_{\mathrm{act}}^sP_{\mathrm{act}}^w\right)
        \notag\\
        &=\sum_j\left(\sum_q\lambda_{j,q}^{|w|+1}\right)
          \tr\!\left(C_j^sC_j^w\right).
        \label{eq:cpsv_full_coordinate_marked_trace_identity}
    \end{align}
    The exponent $|w|+1$ includes the factor at the marked site. Inactive
    coordinates contribute zero.
\end{theorem}

\begin{proof}\leanok
    \uses{thm:cpsv_active_representative_sector_decomposition,
        lem:marked_representative_trace_expansion}
    Split the finite sum over listed coordinates into active and inactive
    parts. The inactive marked blocks vanish. Reindex the active part by the
    class-copy pairs $(j,q)$ and use the bond-dimension identifications to
    replace each active block by its unit-phase representative. The marked
    block supplies one factor $\lambda_{j,q}$ and the word supplies
    $|w|$ further factors. Summing first over $q$ gives
    (\ref{eq:cpsv_full_coordinate_marked_trace_identity}).
\end{proof}

\begin{theorem}[Retained-coordinate representative-grouped Lemma L]
    \label{thm:cpsv_retained_representative_grouped_lemma_l}
    \lean{MPSTensor.CPSVCanonicalFormData.ActiveBNTRefinement.%
        groupedMarkedTensor_basis_eq_of_trace_agree}
    \lean{MPSTensor.CPSVCanonicalFormData.ActiveBNTRefinement.%
        insertedTensor_representative_eq_of_firstSiteActionAgree}
    \leanok
    \uses{thm:cpsv_full_coordinate_marked_trace_identity,
        thm:cpsv_active_representative_sector_decomposition,
        thm:representative_grouped_marked_separation,
        thm:representative_grouped_insert_eq}
    Let $(C_j^s)_j$ and $(E_j^s)_j$ be two marked families on the active
    representatives. Form their full-coordinate marked tensors by retaining
    the inactive listed coordinates as zero blocks. If
    \begin{align}
        \tr\!\left(\widehat C_{\mathrm{grp}}^sT_{\mathrm{grp}}^w\right)
        &=\tr\!\left(\widehat E_{\mathrm{grp}}^sT_{\mathrm{grp}}^w\right)
        \label{eq:cpsv_retained_marked_trace_agreement}
    \end{align}
    for every $s$ and every word $w$, then $C_j=E_j$ for every active
    representative $j$.

    In particular, let $Y,Z\in\MN{d}$. If their first-site actions agree on
    the positive-length matrix product vectors of $T_{\mathrm{grp}}$, then
    \begin{align}
        (YC_j)^i&=(ZC_j)^i
        \label{eq:cpsv_retained_first_site_action}
    \end{align}
    for every active representative $j$ and every physical letter $i$.
\end{theorem}

\begin{proof}\leanok
    \uses{thm:cpsv_full_coordinate_marked_trace_identity,
        thm:cpsv_active_representative_sector_decomposition,
        thm:representative_grouped_marked_separation,
        thm:representative_grouped_insert_eq}
    Theorem~\ref{thm:cpsv_full_coordinate_marked_trace_identity} rewrites
    (\ref{eq:cpsv_retained_marked_trace_agreement}) into marked-trace
    agreement for $P_{\mathrm{act}}$. Its representatives have eventual
    word-tuple separation, and every copy weight is nonzero. Theorem~\ref{thm:representative_grouped_marked_separation}
    therefore gives $C_j=E_j$.

    First-site agreement for $T_{\mathrm{grp}}$ also holds for
    $P_{\mathrm{act}}$ because their matrix product vectors agree at every
    positive length. Applying
    Theorem~\ref{thm:representative_grouped_insert_eq} gives
    (\ref{eq:cpsv_retained_first_site_action}).
\end{proof}

% Source anchors: CPSV16 Appendix C.3, lines 1835--1858, and
% Proposition 4.13, lines 1873--1887.
\begin{theorem}[Original-coordinate marked Lemma L]
    \label{thm:cpsv_original_coordinate_marked_lemma_l}
    \lean{MPSTensor.trace_marked_mul_evalWord_of_coisometry_reconstruction}
    \lean{MPSTensor.linearMarkedTensor_coisometry_reconstruction}
    \lean{MPSTensor.trace_linearMarkedTensor_mul_evalWord_of_coisometry_reconstruction}
    \lean{MPSTensor.CPSVCanonicalFormData.ActiveBNTRefinement.%
        linearMarkedTensor_groupedTensor_eq_groupedMarkedTensor}
    \lean{MPSTensor.CPSVCanonicalFormData.ActiveBNTRefinement.%
        linearMarkedTensor_eq_of_trace_agree}
    \lean{MPSTensor.IsCPSVCanonicalForm.linearMarkedTensor_eq_of_trace_agree}
    \leanok
    \uses{thm:cpsv_canonical_form_active_bnt_refinement,
        thm:cpsv_full_coordinate_marked_trace_identity,
        thm:cpsv_retained_representative_grouped_lemma_l}
    Let $A$ be in literal CPSV canonical form. For canonical-form data and
    its active BNT refinement, write
    \begin{align}
        A^i&=V^\dagger B^iV,
        \notag\\
        VV^\dagger&=\Id.
        \label{eq:cpsv_original_coisometry_reconstruction}
    \end{align}
    where $B^i=GT_{\mathrm{grp}}^iG^{-1}$, $G$ is the listed block gauge,
    and $T_{\mathrm{grp}}$ retains every inactive listed coordinate with
    weight zero. If marked matrices satisfy $C=V^\dagger EV$, then for every
    word $w$, including the empty word,
    \begin{align}
        \tr(CA^w)&=\tr(EB^w).
        \label{eq:cpsv_original_marked_trace}
    \end{align}
    The marked matrix makes the closed chain nonempty; no equality of
    unmarked empty-word matrix product vectors is asserted.

    For a linear physical mark $f$, the linear marking of
    $T_{\mathrm{grp}}$ is exactly its full grouped marked tensor. Each active
    copy $(j,q)$ carries
    $\lambda_{j,q}f(C_j)$, where
    $\lambda_{j,q}=\nu_{k(j,q)}\zeta_{k(j,q)}$, while every inactive marked
    block is zero. For two linear marks $f$ and $g$, put
    $F_A^u=\sum_i f_{ui}A^i$ and $G_A^u=\sum_i g_{ui}A^i$. If
    \begin{align}
        \tr\!\left(F_A^uA^w\right)
        &=\tr\!\left(G_A^uA^w\right)
        \label{eq:cpsv_original_linear_marked_agreement}
    \end{align}
    for every marked letter $u$ and every word $w$, then $F_A=G_A$ on the
    original bond space.
\end{theorem}

\begin{proof}\leanok
    \uses{thm:cpsv_full_coordinate_marked_trace_identity,
        thm:cpsv_retained_representative_grouped_lemma_l}
    Induction on $w$ in (\ref{eq:cpsv_original_coisometry_reconstruction})
    gives
    $CA^w=V^\dagger EB^wV$. Cyclicity of trace and $VV^\dagger=\Id$ give
    (\ref{eq:cpsv_original_marked_trace}). Linearity of the physical mark
    gives the grouped marked tensor, with inactive blocks killed by their
    zero weights.

    Apply (\ref{eq:cpsv_original_marked_trace}) to both sides of
    (\ref{eq:cpsv_original_linear_marked_agreement}), remove the block gauge
    by cyclicity of trace, and use
    Theorem~\ref{thm:cpsv_retained_representative_grouped_lemma_l} on the
    active representatives. The resulting equality of the full grouped
    marked tensors, followed by
    (\ref{eq:cpsv_original_coisometry_reconstruction}), gives $F_A=G_A$.
\end{proof}

\begin{definition}[First-site action on a physical block]
    \label{def:first_site_action_physical_blocking}
    \lean{MPSTensor.firstSiteActionOnBlock}
    \leanok
    \uses{def:blocked_tensor}
    Let $Y\in\MN{d}$, and write physical indices of an $(L+1)$-site block
    as $I=(i_0,\ldots,i_L)$ and $J=(j_0,\ldots,j_L)$. The induced matrix on
    the blocked physical space is
    \begin{align}
        Y^{[L+1]}_{I,J}
        &=Y_{i_0,j_0}\prod_{k=1}^{L}\delta_{i_k,j_k}.
        \label{eq:mpdo_first_site_block}
    \end{align}
\end{definition}

\begin{lemma}[First-site actions under physical blocking]
    \label{lem:first_site_action_physical_blocking}
    \lean{MPSTensor.firstSiteActionAgree_iff_trace}
    \lean{MPSTensor.FirstSiteActionAgree.trace_evalWord}
    \lean{MPSTensor.insertedTensor_firstSiteActionOnBlock_blockTensor}
    \lean{MPSTensor.FirstSiteActionAgree.blockTensor}
    \lean{MPSTensor.insertedTensor_eq_of_firstSiteActionOnBlock_blockTensor_eq}
    \leanok
    \uses{def:first_site_action_physical_blocking, def:blocked_tensor,
        def:l_blk_injective}
    This auxiliary statement supplies the identity needed to combine the
    block-injectivity argument
    \cite[lines~318--345]{Cirac2016MPDO_arXiv} with Appendix~C.3, Lemma~L
    \cite[lines~1835--1858]{Cirac2016MPDO_arXiv}.

    Let $A$ be a tensor, and let $Y^{[L+1]}$ be the physical-block action
    induced by $Y\in\MN{d}$. For $I=(i_0,\ldots,i_L)$, its insertion into
    the blocked tensor satisfies
    \begin{align}
        (Y^{[L+1]}A^{[L+1]})^I
        &=(YA)^{i_0}A^{i_1}\cdots A^{i_L}.
        \label{eq:mpdo_blocked_insertion}
    \end{align}
    Hence equality of the first-site actions of $Y,Z$ on the MPV family of
    $A$ implies equality of the first-block actions of
    $Y^{[L+1]},Z^{[L+1]}$ on the MPV family of $A^{[L+1]}$.

    The first-site actions of $Y,Z$ agree at every length if and only if
    \begin{align}
        \tr((YA)^i A^w)
        &=\tr((ZA)^i A^w),
        \label{eq:mpdo_marked_first_site_trace}
    \end{align}
    for every physical index $i$ and every finite word $w$.

    Conversely, suppose that $A$ is $L$-block injective. If the two induced
    insertion tensors on $A^{[L+1]}$ are equal, then $YA=ZA$.
\end{lemma}

\begin{proof}
    \leanok
    For the forward implication, split the first block into its first letter
    and its remaining $L$ letters. The insertion identity
    (\ref{eq:mpdo_blocked_insertion}) reduces the blocked MPV equality to the
    original first-site equality, with the remaining letters of the first
    block prepended to the flattened tail.

    For the converse, put $\Delta^i=(YA)^i-(ZA)^i$. Equality after blocking
    gives $\Delta^i A^w=0$ for every word $w$ of length $L$.
    Since the length-$L$ words span $\MN{D}$, the same equality holds with
    $A^w$ replaced by the identity, and therefore $\Delta^i=0$ for every
    physical index $i$.
\end{proof}

\begin{theorem}[Representative separation after injective blocking]
    \label{thm:postblocked_representative_grouped_insert_eq}
    \lean{MPSTensor.IsBNTCanonicalForm.eventuallyRepresentativeWordTupleSpan_of_basis_injective}
    \lean{MPSTensor.IsBNTCanonicalForm.eventuallyRepresentativeWordTupleSpan_blockTensor}
    \lean{MPSTensor.IsBNTCanonicalForm.eventuallyRepresentativeWordTupleSpan}
    \lean{MPSTensor.pairTraceSeparatingAt_blockTensor}
    \lean{MPSTensor.IsBNTCanonicalForm.insertedTensor_basis_eq_of_firstSiteActionAgree_of_basis_injective}
    \lean{MPSTensor.IsBNTCanonicalForm.insertedTensor_basis_eq_of_sameMPV₂Pos_firstSiteActionAgree_of_basis_injective}
    \lean{MPSTensor.IsBNTCanonicalForm.insertedTensor_basis_eq_of_firstSiteActionAgree_of_common_blockInjective}
    \lean{MPSTensor.IsBNTCanonicalForm.insertedTensor_basis_eq_of_sameMPV₂Pos_firstSiteActionAgree_of_common_blockInjective}
    \lean{MPSTensor.IsBNTCanonicalForm.exists_common_basis_isNBlkInjective}
    \lean{MPSTensor.IsBNTCanonicalForm.insertedTensor_basis_eq_of_firstSiteActionAgree}
    \lean{MPSTensor.IsBNTCanonicalForm.insertedTensor_basis_eq_of_sameMPV₂Pos_firstSiteActionAgree}
    \leanok
    \uses{def:sector_bnt_cf, def:sector_decomposition_block_tensor,
      def:representative_word_tuple_span, thm:representative_grouped_insert_eq}
    Let $P$ be a BNT canonical form with basis representatives
    $(A_j)_{j=0}^{g-1}$, and suppose that the representatives have already
    been blocked so that every $A_j$ is injective. Then representative
    word-tuple separation holds at every length $L\geq 6(g-1)+1$.
    Consequently, if the first-site actions of $Y,Z\in\MN{d}$ agree on the
    positive-length MPV family of the tensor represented by $P$, then
    $YA_j=ZA_j$ for $0\leq j<g$.
    The same conclusion holds when the first-site equality is given for any
    tensor with the same positive-length MPV family as the tensor represented by
    $P$.

    More generally, let $P$ be a BNT canonical form before physical blocking,
    let $p>0$, and suppose that every $p$-blocked representative
    $A_j^{[p]}$ is injective. Then the blocked sector decomposition $P^{[p]}$
    has representative word-tuple separation at every blocked length
    $L\geq 6(g-1)+1$.

    Without changing the physical alphabet, a BNT canonical form also has
    eventual representative word-tuple separation. If $p>0$ is a common
    block-injectivity length, then one may take pairwise separating words of
    length $6p$, form selectors of length $6p(g-1)$, and append homogeneous
    words of any length at least $p$.

    Finally, the BNT canonical-form hypotheses imply that there is a common
    $L>0$ for which every unblocked representative $A_j$ is $L$-block
    injective. Indeed, irreducibility and left-canonicality give a finite
    peripheral period $m_j$. The self-overlap tends to $m_j$ along multiples
    of $m_j$, but by normalization it also tends to $1$; hence $m_j=1$.
    Thus every representative is normal, and finiteness gives a common
    positive block-injectivity length. If the first-site actions
    of $Y,Z$ agree on the tensor represented by $P$, then
    $YA_j=ZA_j$ for $0\leq j<g$.
    The same conclusion holds when the first-site equality is given for any
    tensor with the same positive-length MPV family. This composes the blocking
    conclusion of \cite[lines~318--345]{Cirac2016MPDO_arXiv} with
    Appendix~C.3, Lemma~L
    \cite[lines~1835--1858]{Cirac2016MPDO_arXiv}.

    No additional common-length hypothesis is required.
    % The remaining MPDO scope question is the relation with the horizontal
    % copy-indexed hypotheses; see docs/paper-gaps/cpgsv17_bicf_block_separation.tex.
\end{theorem}

\begin{proof}
    \leanok
    \uses{lem:bnt_direct_sum_block_separation_word_span,
      lem:block_separating_equations_from_pairwise_separation,
      lem:product_word_span_from_bnt_block_separation,
      thm:representative_grouped_insert_eq}
    Injectivity of $A_j$ implies that its homogeneous words span
    $\MN{D_j}$ at every positive length. Apply the three-block separation
    argument with source length $1$. For each ordered pair of distinct
    representatives it gives a pairwise block-separating equation at length
    $6$. Multiplying the equations for the $g-1$ representatives different
    from $A_j$ gives coefficients $c_w^{(j)}$ satisfying
    \begin{align}
        \sum_{|w|=6(g-1)}c_w^{(j)}A_k^w
        &=
        \begin{cases}
            \Id, & k=j,\\
            0, & k\neq j.
        \end{cases}.
        \label{eq:mpdo_rep_selectors}
    \end{align}
    If $L\geq6(g-1)+1$, put $m=L-6(g-1)>0$. The length-$m$ words of each
    $A_j$ span $\MN{D_j}$. Concatenating these words with the selectors
    in (\ref{eq:mpdo_rep_selectors}) shows that the simultaneous length-$L$
    word tuples span $\prod_j\MN{D_j}$. This proves eventual representative
    word-tuple separation. The representative-grouped Lemma L gives the two
    first-site conclusions.

    For the unblocked canonical form, propagate $p$-block injectivity to the
    lengths $2p-1$, $2p$, and $6p$. Applying the direct-sum argument to the
    resulting injective tensors gives pair trace separation at original
    length $6p$. The
    bijection between blocked words of length $6$ and original words of length
    $6p$ gives pair trace separation for the $p$-blocked
    representatives. The preceding selector argument then applies verbatim
    to $P^{[p]}$.

    The same pairwise equations at original length $6p$ give selectors for
    the unblocked representatives directly. For every
    $L\geq6p(g-1)+p$, the remaining prefix has length at least $p$ and hence
    spans each full representative matrix algebra. Concatenating the prefix
    with the selectors proves eventual representative word-tuple separation
    before blocking.

    For the final assertion, first consider one representative $A_j$.
    Irreducibility and left-canonicality give a positive peripheral period
    $m_j$. The periodic self-overlap formula and the normalized self-overlap
    hypothesis give, along the same subsequence,
    \begin{align}
        \lim_{k\to\infty}O_{A_jA_j}(m_jk)
        &=m_j,
        \notag\\
        \lim_{k\to\infty}O_{A_jA_j}(m_jk)
        &=1.
        \notag
    \end{align}
    Hence $m_j=1$.
    Primitivity then implies normality. Taking a common
    multiple of the finitely many normality lengths gives a positive length
    $L$ at which every representative is block injective. Block $L+1$ sites.
    Injectivity
    propagates from length $L$ to length $L+1$, so the blocked representatives
    are injective. Act on the first letter of each physical block and apply
    the blocked representative separation followed by the
    representative-grouped Lemma~L. This gives equality of the two induced
    insertions on every blocked representative. The remaining $L$ letters
    of the first block form a spanning family, so
    Lemma~\ref{lem:first_site_action_physical_blocking} recovers
    $YA_j=ZA_j$ before blocking.
\end{proof}

\begin{corollary}[Marked-letter separation for a BNT canonical form]
    \label{cor:bnt_marked_representative_separation}
    \lean{MPSTensor.IsBNTCanonicalForm.markedTensor_basis_eq_of_trace_agree}
    \leanok
    \uses{def:sector_bnt_cf,
        thm:postblocked_representative_grouped_insert_eq,
        thm:representative_grouped_marked_separation}
    Let $(C_j^s)$ and $(E_j^s)$ be two marked families over a BNT canonical
    form. If their closed marked chains agree for every marked index $s$ and
    every unmarked word $w$, then $C_j^s=E_j^s$ for every representative $j$
    and every marked index $s$.
\end{corollary}

\begin{proof}
    \leanok
    \uses{thm:postblocked_representative_grouped_insert_eq,
        thm:representative_grouped_marked_separation}
    The BNT canonical-form hypotheses give eventual representative word-tuple
    separation. Apply
    Theorem~\ref{thm:representative_grouped_marked_separation}.
\end{proof}

\begin{corollary}[Physical-letter marked separation in BNT-refined horizontal form]
    \label{cor:horizontal_cf_linear_marked_separation}
    \lean{MPSTensor.linearMarkedTensor}
    \lean{MPSTensor.linearMarkedTensor_toTensorFromBlocks}
    \lean{MPSTensor.SectorDecomposition.linearMarkedTensor_toTensor_eq_markedTensor}
    \lean{MPSTensor.linearMarkedTensor_gauge}
    \lean{MPSTensor.trace_linearMarkedTensor_mul_evalWord_gauge}
    \lean{MPOTensor.IsHorizontalCF.linearMarkedTensor_eq_of_trace_agree}
    \leanok
    \uses{def:horizontal_cf_mpdo,
        cor:bnt_marked_representative_separation}
    Let $M$ be an MPO tensor in normalized BNT-refined horizontal form. For two
    coefficient families $f_{u,z}$ and $g_{u,z}$, put
    $C^u=\sum_z f_{u,z}M^z$ and $E^u=\sum_z g_{u,z}M^z$.
    If $\tr(C^uM^w)=\tr(E^uM^w)$ for every marked index $u$ and every word
    $w$, then $C^u=E^u$ for every $u$.

    The marks are required to lie in the span of the physical letters. No
    assertion is made for arbitrary bond-space matrices, since off-diagonal
    matrices between repeated sectors are invisible to closed chains. The
    corollary is the physical-insertion consequence of
    \cite[Appendix~C.3, lines~1835--1858]{Cirac2016MPDO_arXiv} used in the
    contractions of
    \cite[Proposition~4.13, lines~1909--1919]{Cirac2016MPDO_arXiv}.
\end{corollary}

\begin{proof}
    \leanok
    \uses{cor:bnt_marked_representative_separation}
    Pass through the block-diagonal similarity in the normalized BNT-refined
    horizontal form. Linear combinations of physical letters are covariant
    under this similarity, and cyclicity of the trace removes it from each
    closed marked chain. On the representative-indexed sector
    decomposition one has
    \begin{align}
        C_{\mathrm{tot}}^u
        &=\bigoplus_{j,q}\mu_{j,q}
          \sum_z f_{u,z}A_j^z,
        \notag\\
        E_{\mathrm{tot}}^u
        &=\bigoplus_{j,q}\mu_{j,q}
          \sum_z g_{u,z}A_j^z.
        \notag
    \end{align}
    Corollary~\ref{cor:bnt_marked_representative_separation} identifies the
    two marks on every minimal representative. Taking their weighted direct
    sums and restoring the common block-diagonal similarity gives
    $C^u=E^u$.
\end{proof}
